\documentclass[a4paper,10pt]{article}

% --- PREÂMBULO ---
\usepackage[a4paper, top=0.9cm, bottom=0.8cm, left=1.5cm, right=1.5cm]{geometry}
\usepackage[T1]{fontenc}
\usepackage[utf8]{inputenc}
\usepackage[provide=*,portuguese]{babel}

% O template usa fonte serifada (Times-like); mantemos a família padrão
% do LaTeX substituída por Times para fidelidade visual ao modelo.
\usepackage{mathptmx}

\usepackage{enumitem}
\setlist[itemize]{
    label=\textbullet,
    leftmargin=1.5em,
    itemsep=1pt,
    topsep=2pt,
    parsep=0pt
}

\usepackage{hyperref}
\usepackage{titlesec}
\usepackage[svgnames]{xcolor}
\usepackage{amsmath}
\usepackage{amssymb}

\definecolor{linkcolor}{RGB}{0, 0, 238}

\hypersetup{
    colorlinks=true,
    urlcolor=linkcolor,
    linkcolor=linkcolor,
    pdftitle={Israel Souza Ferreira - CV},
    pdfauthor={Israel Souza Ferreira},
}

% --- SEÇÕES NO ESTILO DO TEMPLATE ---
% Título em negrito, caixa alta, alinhado à esquerda, com régua fina abaixo.
\titleformat{\section}
  {\normalsize\bfseries}
  {}
  {0em}
  {\MakeUppercase}
  [\vspace{-1.2ex}\rule{\linewidth}{0.4pt}]

\titlespacing{\section}{0pt}{1.6ex}{0.8ex}

% --- MACRO DE ENTRADA (Empresa/Universidade + local + cargo + datas) ---
% #1 = Empresa (negrito, esq.)   #2 = (regime) - país (negrito, dir.)
% #3 = Cargo (itálico, esq.)     #4 = Datas (itálico, dir.)
\newcommand{\entrada}[4]{%
    \noindent\textbf{#1} \hfill \textbf{#2}\\
    \textit{#3} \hfill \textit{#4}\par
}

\pagestyle{empty}
\raggedbottom
\linespread{0.98}

\begin{document}

% ---------- CABEÇALHO ----------
\begin{center}
    {\LARGE \textbf{Israel Souza Ferreira}} \\[3pt]

    {\small
    Fortaleza, Ceará, Brasil (Disponível para trabalho remoto)
    \; | \;
    +55 (85) 984332790
    \; | \;
    \href{mailto:israel.souza.dev@gmail.com}{e-mail}
    \; | \;
    \href{https://www.linkedin.com/in/isrreal/}{LinkedIn}
    \; | \;
    \href{https://github.com/isrreal}{GitHub}
    }
\end{center}

% ---------- RESUMO PROFISSIONAL ----------
\section*{Resumo Profissional}

Pesquisador de \textit{Machine Learning} e graduando em Ciência da Computação
pela Universidade Federal do Ceará (conclusão em julho de 2026). Combino
rigor matemático e profundidade de engenharia em PyTorch, C++, sistemas
multimodais e MLOps. Desenvolvi um regressor multitarefa multimodal para
2.580 questões do ENEM, estimando simultaneamente parâmetros da TRI, curva
ICC, taxa de acerto e classificações pedagógicas. Possuo experiência no ciclo
completo de pesquisa aplicada e engenharia de ML, incluindo a primeira
abordagem meta-heurística da literatura para um problema NP-completo em
grafos e o desenvolvimento solo de um produto de visão computacional em
fase final de homologação para um cliente real.

% ---------- EXPERIÊNCIA PROFISSIONAL ----------
\section*{Experiência Profissional}

\entrada
    {Tieta Artificial Intelligence}
    {(remoto) -- Brasil}
    {Pesquisador de Machine Learning (Bolsista CNPq RHAE)}
    {Abr 2025 -- Presente}

\begin{itemize}

    \item Atuei em pesquisa aplicada de \textit{Machine Learning} e
    engenharia de dados, abrangendo um sistema de recuperação multimodal e
    um regressor multitarefa para análise de questões do ENEM; realocado em
    jul/2026 para novo projeto de pesquisa com continuidade até dez/2026.

    \item Desenvolvi um \textbf{regressor multitarefa multimodal} para
    2.580 questões do ENEM, executando simultaneamente seis tarefas:
    predição dos parâmetros da TRI $(a,b,c)$, reconstrução da curva ICC,
    estimativa da taxa de acerto e classificação de área, competência e
    habilidade pedagógica.

    \item Modelei a arquitetura com princípios de \textbf{MoE/MMoE}
    (\textit{experts}, redes de \textit{gating}, atenção e
    \textit{cross-attention}) e implementei redutores vetoriais baseados em
    LSTM e \textit{Attention Heads} para condensar representações
    multimodais antes das cabeças de cada tarefa.

    \item Estruturei a entrada multimodal com \textit{embeddings} textuais do
    Portulan Serafim-900M e representações visuais compactadas pelo DCAE,
    avaliando estratégias de redução e fusão texto--imagem.

    \item Concebi um canal de \textbf{raciocínio supervisionado} com DSPy
    (\texttt{Predict} + \texttt{ChainOfThought}), no qual LLMs locais
    resolvem e justificam as questões; as cadeias geradas alimentam o
    regressor como representação adicional, com ganho consistente validado
    pelo teste de Wilcoxon na estimativa de dificuldade latente.

    \item Conduzi \textit{feature engineering} sobre a geometria das
    trajetórias inferenciais (comprimento, curvatura, linearidade),
    diagnostiquei a incompatibilidade do \textit{zero-padding} no espaço
    RelRep, propus injeção lateral via MLP e demonstrei por ablação a
    redundância dessas métricas frente ao canal de raciocínio.

    \item Desenvolvi uma \textbf{API de busca multimodal top-$k$} em FastAPI
    com representações relativas (RelRep), similaridade BERTScore, fusão
    \textit{early}/\textit{late} e indexação vetorial em FAISS para
    recuperação e ranqueamento de questões.

    \item Formulei um arcabouço de métricas geométricas para seleção de
    âncoras (cobertura, posto efetivo, número de condição,
    \textit{kNN-overlap}), comparando \textit{k-center greedy}, LogDet e
    seleção aleatória, com escolha por dominância de Pareto e validação por
    Recall@K.

    \item Implementei o pipeline de treinamento em PyTorch (Focal Loss, BCE,
    MSE, Adam/AdamW, Cosine Annealing) e análises estatísticas (Wilcoxon,
    Qui-quadrado, McNemar, Pearson/Spearman, ROC, KDE), com experimentos
    reprodutíveis via Docker, Git e GitLab CI/CD.

\end{itemize}

\vspace{4pt}

\entrada
    {Cliente autônomo (instituição de ensino)}
    {(remoto) -- Brasil}
    {Desenvolvedor de Machine Learning --- Ponto com Biometria Facial}
    {2026 -- Presente}

\begin{itemize}

    \item Desenvolvi \textbf{sozinho} um sistema completo de registro de
    ponto por reconhecimento facial --- modelagem de ML, \textit{backend},
    segurança, banco de dados e \textit{deploy} --- atualmente em fase final
    de homologação para o cliente.

    \item Implementei reconhecimento facial 1:1 e 1:N com
    \textit{embeddings} InsightFace (\texttt{buffalo\_l}) em
    $\mathbb{R}^{512}$, verificação por similaridade de cosseno e
    identificação por KNN executada no PostgreSQL via pgvector; robustez por
    agregação \textit{multi-frame}, \textit{liveness} e proteção anti-replay
    com SHA-256.

    \item Treinei um classificador de atestados EfficientNet-B0
    ($\sim$85\% de acurácia) com \textit{fine-tuning} seletivo, Optuna/TPE,
    Stratified K-Fold e AMP; decisão híbrida CNN + OCR isolada em
    microsserviço de inferência.

    \item Construí \textit{backend} FastAPI com concorrência sync/async por
    perfil de carga, JWT com revogação persistida, \textit{geofencing}
    (Haversine + perímetro elíptico) e \textit{hardening} de segurança
    (PIN de provisionamento, comparação em tempo constante, anti-IDOR);
    persistência com SQLAlchemy, Alembic e PostgreSQL.

    \item Estruturei MLOps de ponta a ponta: realimentação contínua de
    re-treino por concordância humano-modelo, entrega Dockerizada com HTTPS
    via Cloudflare Tunnel, 115 testes automatizados e integração contínua no
    GitHub Actions.

\end{itemize}

\vspace{4pt}

\entrada
    {Universidade Federal do Ceará --- Campus Quixadá}
    {(presencial) -- Brasil}
    {Monitor de Algoritmos, Cálculo e Pré-Cálculo}
    {2022 -- 2024}

\begin{itemize}
    \item Ministrei sessões de apoio em estruturas de dados, algoritmos e
    análise de complexidade, além de Cálculo Diferencial e Integral aplicado
    à Computação.
\end{itemize}

% ---------- HABILIDADES ----------
\section*{Habilidades}

\noindent
\textbf{Linguagens:} Python, C++, SQL, Bash.\\[2pt]
\textbf{Frameworks \& Bibliotecas:} PyTorch, Scikit-learn, DSPy, FAISS,
Optuna, InsightFace, OpenCV, FastAPI, SQLAlchemy, Alembic, Pandas, NumPy,
SciPy, Matplotlib.\\[2pt]
\textbf{Ferramentas \& Plataformas:} Docker, Git, GitLab CI/CD,
GitHub Actions, PostgreSQL, pgvector, MLflow.\\[2pt]
\textbf{Arquiteturas \& Métodos:} MMoE, MoE, PLE, LSTM, atenção e
\textit{cross-attention}, RelRep, BERTScore, busca vetorial top-$k$,
testes de hipótese, TRI/ICC, ROC, KDE, SVD, matriz de Gram, Nyström.

% ---------- IDIOMAS ----------
\section*{Idiomas}

\noindent
\textbf{Português:} Nativo \; | \;
\textbf{Inglês:} Leitura técnica fluente (artigos científicos e
documentações)

% ---------- FORMAÇÃO ----------
\section*{Formação Acadêmica}

\entrada
    {Universidade Federal do Ceará --- Campus Quixadá}
    {Brasil}
    {Bacharelado em Ciência da Computação}
    {2021 -- Jul 2026}

% ---------- DESENVOLVIMENTO PROFISSIONAL ----------
\section*{Desenvolvimento Profissional (Projetos e Atividades)}

\begin{itemize}

    \item \textbf{\href{https://github.com/isrreal/Triple-Roman-Domination-in-graphs}
    {3RDF --- Triple Roman Domination Problem (TCC, 2024--2025):}}
    primeira abordagem meta-heurística da literatura para este problema
    NP-completo em grafos (C++), com Algoritmo Genético, ACO
    (MAX-MIN Ant System + RVNS) e formulação de Programação Linear Inteira
    como \textit{baseline} ótimo de validação experimental.

    \item \textbf{\href{https://github.com/isrreal/Dengue-Forecasting-System}
    {Dengue Forecasting System (2025):}}
    sistema MLOps em desenvolvimento para previsão de notificações de dengue
    por estado; ETL do SINAN via pySUS, PostgreSQL, MLflow,
    FastAPI, \textit{dashboard} Streamlit/Plotly, pytest, mypy e
    Docker Compose.

\end{itemize}

\end{document}
